% !TeX root = main.tex
\documentclass[manuscript,review,nonacm]{acmart}

% nonacm suppresses publication metadata for this working draft.
% Remove nonacm and insert the venue's supplied rights/conference metadata
% when preparing an ACM publication. For an anonymous review submission,
% follow the venue's instructions (often anonymous,review options).
\settopmatter{printacmref=false}
\setcopyright{none}
\renewcommand\footnotetextcopyrightpermission[1]{}

\usepackage{booktabs}
\usepackage{mathpartir}

\newcommand{\G}{\mathsf{G}}
\newcommand{\E}{\mathsf{E}}
\newcommand{\real}[1]{\mathsf{real}^{#1}}
\newcommand{\bool}{\mathsf{bool}}
\newcommand{\judg}[4]{#1 \mid #2 \vdash #3 : #4} 
\newcommand{\sep}{,}
\newcommand{\emptybunch}{\varnothing}
\newcommand{\hole}[2]{#1[#2]}
\newcommand{\wf}{\mathsf{wf}}

\title{[Project Title]: Extension}
\author{[Author Name]}
\affiliation{%
  \institution{[Institution]}
  \city{[City]}
  \country{[Country]}}
\email{author@example.com}

\begin{document}

\maketitle

\section{NEW}
\paragraph{Types}
$\tau$ describes all expression types.

\begin{align*}
\tau ::= {}&
\real{\G} \mid \real{\E}
\mid \mathsf{bool}
\mid \mathsf{unit}
\mid \tau \times \tau
\mid \tau \otimes \tau.\\
\sigma ::= {}&
\real{\G} \mid \mathsf{bool} \mid \mathsf{unit}
\mid \sigma \times \sigma
\end{align*}

\paragraph{Typing contexts.}
Typing judgments have the form $\Delta \mid \Gamma \vdash e : \tau$, where $\Delta$ contains variables of type $\sigma$ that may be used freely across subexpressions, and $\Gamma$ tracks dependencies of values that must be separated when a typing rule requires independence.
\begin{align*}
\Delta &::= \varnothing \mid \Delta,x:\sigma \\
\Gamma &::= \varnothing
  \mid x:\tau
  \mid (\Gamma;\Gamma)
  \mid (\Gamma,\Gamma)
\end{align*}

\begin{enumerate}
  \item $\Gamma_1 ; \Gamma_2$ permits sharing (additive).
  \item The $\Gamma_1 , \Gamma_2$ reqires separation (multiplicative).
\end{enumerate}

Here, $\Gamma$ is a bunched (tree-like) context, whereas $\Delta$ is an ordinary context.

\paragraph{One-hole bunched contexts.}
We work with a context containing a single hole that may be filled by another context. We use the metavariable $\Gamma_0$ to range over contexts with a hole and write $\Gamma_0(\Gamma)$ to fill the hole in $\Gamma_0$ with $\Delta$.
\begin{align*}
\Gamma_0 ::= {}&
\cdot
\mid (\Gamma_0;\Gamma)
\mid (\Gamma;\Gamma_0)
\mid (\Gamma_0,\Gamma)
\mid (\Gamma,\Gamma_0).
\end{align*}
where $\cdot$ denotes the unique hole here.

For example,
\[
\Gamma_0=(\cdot,z:\real{\E})
\quad\Longrightarrow\quad
\Gamma_0(x:\real{\E};y:\real{\E})
=((x:\real{\E};y:\real{\E}),z:\real{\E}).
\]

\paragraph{Well-formed bunches}
NOTE:*** Additive sharing should permit two occurrences of the same variable, but not with contradictory types.

A bunch $\Gamma$ is well-formed, written $\wf(\Gamma)$, if shared variables have consistent types and no variable occurs on both sides of a multiplicative comma.

Formally, we can define this as:
\begin{align*}
\wf(\varnothing) &\iff \top,\\
\wf(x:\tau) &\iff \top,\\
\wf(\Gamma_1;\Gamma_2)
&\iff \wf(\Gamma_1)\land\wf(\Gamma_2) \land
\forall x\in\operatorname{dom}(\Gamma_1)\cap
\operatorname{dom}(\Gamma_2),\
\Gamma_1(x)=\Gamma_2(x),\\
\wf(\Gamma_1,\Gamma_2)
&\iff \wf(\Gamma_1)\land\wf(\Gamma_2) \land
\operatorname{dom}(\Gamma_1)\cap
\operatorname{dom}(\Gamma_2)=\varnothing.
\end{align*}

where $\Gamma(x)$ denotes the type assigned to $x$ in $\Gamma$.

For example, a well-formed bunch is $(x:\tau;y:\tau),z:\tau$, and a bunch that is not well-formed is $(x:\tau;y:\tau),x:\tau$.



\paragraph{Typing rules.}
Let $m\in\{\G,\E\}$.

\begin{mathpar}
% Variables and constants
\inferrule[Var-$\Delta$]
  { x:\tau \in \Delta }
  { \Delta\mid\varnothing\vdash x:\tau }

\and
\inferrule[Var-$\Gamma$]
  { }
  { \Delta\mid x:\tau
      \vdash x:\tau }

\and
\inferrule[Real]
  { r\in\mathbb R }
  { \Delta\mid\varnothing\vdash r:\real{\G} }

\and
\inferrule[Bool]
  { b\in\{\mathsf{true},\mathsf{false}\} }
  { \Delta\mid\varnothing\vdash b:\bool }

\and
\inferrule[Unit]
  { }
  { \Delta\mid\varnothing\vdash ():\mathsf{unit} }
\end{mathpar}

\begin{mathpar}
% Arithmetic
\inferrule[Add]
  { \Delta\mid\Gamma\vdash e_1:\real{m} \\
    \Delta\mid\Gamma\vdash e_2:\real{m} }
  { \Delta\mid\Gamma
      \vdash e_1+e_2:\real{m} }

\and
\inferrule[Mul-Ind]
  { \Delta\mid\Gamma_1\vdash e_1:\real{\E} \\
    \Delta\mid\Gamma_2\vdash e_2:\real{\E} \\
    \wf(\Gamma_1,\Gamma_2) }
  { \Delta\mid(\Gamma_1,\Gamma_2)
      \vdash e_1\cdot e_2:\real{\E} }

\and
\inferrule[Mul-G]
  { \Delta\mid\varnothing
      \vdash e_1:\real{\G} \\
    \Delta\mid\Gamma\vdash e_2:\real{m} }
  { \Delta\mid\Gamma
      \vdash e_1\cdot e_2:\real{m} }
\end{mathpar}

\begin{mathpar}
% Let-bindings
\inferrule[Let-E]
  { \Delta\mid\Gamma_1\vdash e_1:\tau_1 \\
    \Delta\mid\Gamma_0(x:\tau_1)
      \vdash e_2:\tau_2 \\
    \wf(\Gamma_0(\Gamma_1)) }
  { \Delta\mid\Gamma_0(\Gamma_1)
      \vdash
      \mathsf{let}\ x=e_1\ \mathsf{in}\ e_2:\tau_2 }

\and
\inferrule[Let-G]
  { \Delta\mid\varnothing\vdash e_1:\sigma \\
    \Delta,x:\sigma\mid\Gamma\vdash e_2:\tau_2 }
  { \Delta\mid\Gamma\vdash
      \mathsf{let}\ x=e_1\ \mathsf{in}\ e_2:\tau_2 }
\end{mathpar}

\begin{mathpar}
% Conditionals
\inferrule[If]
  { \Delta\mid\varnothing\vdash e_0:\bool \\
    \Delta\mid\Gamma\vdash e_1:\tau \\
    \Delta\mid\Gamma\vdash e_2:\tau }
  { \Delta\mid\Gamma\vdash
      \mathsf{if}\ e_0\ \mathsf{then}\ e_1\
      \mathsf{else}\ e_2:\tau }
\end{mathpar}

\begin{mathpar}
% Uniform distributions
\inferrule[Uniform-E]
  { \Delta\mid\Gamma\vdash e_1:\real{\E} \\
    \Delta\mid\Gamma\vdash e_2:\real{\E} }
  { \Delta\mid\Gamma\vdash
      \mathsf{uniform}^{\E}(e_1,e_2)
      :\real{\E} }

\and
\inferrule[Uniform-G]
  { \Delta\mid\varnothing
      \vdash e_1:\real{\G} \\
    \Delta\mid\varnothing
      \vdash e_2:\real{\G} }
  { \Delta\mid\varnothing\vdash
      \mathsf{uniform}^{\G}(e_1,e_2)
      :\real{\G} }

\and
\inferrule[Uniform-M]
  { \Delta\mid\Gamma\vdash e_1:\real{m} \\
    \Delta\mid\Gamma\vdash e_2:\real{m} }
  { \Delta\mid\Gamma\vdash
      \mathsf{uniform}^{\mathsf{M}}(e_1,e_2)
      :\real{m} }

\and
\inferrule[Flip-G]
  { \Delta\mid\varnothing\vdash p:\real{\G} }
  { \Delta\mid\varnothing\vdash
      \mathsf{flip}^{\G}(p):\bool }
\end{mathpar}

\paragraph{Structural rules.}
The following rules show weakening, contraction,
and conversion in bunched contexts.

\begin{mathpar}
\inferrule[Weakening]
  { \Delta\mid\Gamma_0(\Gamma)\vdash e:\tau \\
    \wf(\Gamma_0(\Gamma;\Gamma')) }
  { \Delta\mid\Gamma_0(\Gamma;\Gamma')
      \vdash e:\tau }

\and
\inferrule[Contraction]
  { \Delta\mid\Gamma_0(\Gamma;\Gamma)
      \vdash e:\tau }
  { \Delta\mid\Gamma_0(\Gamma)\vdash e:\tau }

\and
\inferrule[Conversion]
  { \Delta\mid\Gamma_0(\Gamma_1;\Gamma_2)
      \vdash e:\tau \\
    \wf(\Gamma_0(\Gamma_1,\Gamma_2)) }
  { \Delta\mid\Gamma_0(\Gamma_1,\Gamma_2)
      \vdash e:\tau }
\end{mathpar}

The \textsc{Weakening} and \textsc{Contraction} rules follow standard BI. The \textsc{Conversion} rule allows conversion from additive to mutiplicative composition when the resulting bunch is well-formed.

\paragraph{Examples}
\begin{example}[Sharing versus separation]
Consider $x+x$ and $x\cdot x$, where $x:\real{\E}$.
By \textsc{Var-$\Gamma$}, both operands have the judgment $\varnothing\mid x:\real{\E}\vdash x:\real{\E}$.

For addition, \textsc{Add} uses the same context $x:\real{\E}$
in both premises and directly gives
$\varnothing\mid x:\real{\E}\vdash x+x:\real{\E}$. For multiplication, \textsc{Mul-Ind} requires: $\wf(x:\real{\E},x:\real{\E})$, but this fails because $x$ occurs on both sides of the comma.
\end{example}

\begin{example}
Consider
\[
\mathsf{let}\ x=u+v\ \mathsf{in}\ x\cdot w,
\]
where $u,v,w$ are distinct variables of type $\real{\E}$.
By \textsc{Weakening}, both $u$ and $v$ can be typed under
the shared context below; \textsc{Add} then types $u+v$ under it:
\[
\Gamma_1=(u:\real{\E};v:\real{\E}).
\]
The semicolon permits $u$ and $v$ to share dependencies.
By \textsc{Mul-Ind}, the body $x\cdot w$ is typed under
\[
x:\real{\E},w:\real{\E}.
\]
The comma requires separation between $x$ and $w$.

To apply \textsc{Let-E}, take
$\Gamma_0=(\cdot,w:\real{\E})$.
Replacing the binding for $x$ with $\Gamma_1$ yields
\[
\varnothing
\mid ((u:\real{\E};v:\real{\E}),w:\real{\E})
\vdash
\mathsf{let}\ x=u+v\ \mathsf{in}\ x\cdot w
:\real{\E}.
\]
\end{example}

\paragraph{Typing rules (for pairs and lists).}
PROBLEM I was having:
\[
\mathsf{let}\ p=(x,y)\ \mathsf{in}\
(\mathsf{fst}\ p)\cdot(\mathsf{snd}\ p).
\]

Even if $x$ and $y$ are independent, the ordinary pair
type $p:\real{\E}\times\real{\E}$ does not record this.

\[
(p:\real{\E}\times\real{\E}),
(p:\real{\E}\times\real{\E}),
\]

Both projections use the same binding $p$.
Applying \textsc{Mul-Ind} requires:
\[
\wf(p:\real{\E}\times\real{\E},
      p:\real{\E}\times\real{\E}),
\]
This fails because $p$ occurs on both sides of the comma.

NEED to preserve independence in the pair's type.


Ordinary pairs $\tau_1\times\tau_2$ may have dependent
components, and independent pairs $\tau_1\otimes\tau_2$
have independent components.

\begin{mathpar}
\inferrule[$\times$-Intro]
  { \Delta\mid\Gamma\vdash e_1:\tau_1 \\
    \Delta\mid\Gamma\vdash e_2:\tau_2 }
  { \Delta\mid\Gamma
      \vdash(e_1,e_2):\tau_1\times\tau_2 }

\and
\inferrule[$\times$-Elim$_1$]
  { \Delta\mid\Gamma\vdash e:\tau_1\times\tau_2 }
  { \Delta\mid\Gamma\vdash\mathsf{fst}\ e:\tau_1 }

\and
\inferrule[$\times$-Elim$_2$]
  { \Delta\mid\Gamma\vdash e:\tau_1\times\tau_2 }
  { \Delta\mid\Gamma\vdash\mathsf{snd}\ e:\tau_2 }
\end{mathpar}

\begin{mathpar}
\inferrule[$\otimes$-Intro]
  { \Delta\mid\Gamma_1\vdash e_1:\tau_1 \\
    \Delta\mid\Gamma_2\vdash e_2:\tau_2 \\
    \wf(\Gamma_1,\Gamma_2) }
  { \Delta\mid(\Gamma_1,\Gamma_2)
      \vdash (e_1, e_2):\tau_1\otimes\tau_2 }

\and
\inferrule[$\otimes$-Elim]
  { \Delta\mid\Gamma_1\vdash e_1:\tau_1\otimes\tau_2 \\
    \Delta\mid\Gamma_0(x:\tau_1,\,y:\tau_2)\vdash e_2:\tau \\
    \wf(\Gamma_0(\Gamma_1)) }
  { \Delta\mid\Gamma_0(\Gamma_1)\vdash
    \mathsf{let}\ (x,y)=e_1\ \mathsf{in}\ e_2:\tau }
\end{mathpar}

\paragraph{Examples}
\begin{example}
Consider
\[
\mathsf{let}\ (a,b)=(x,y)\ \mathsf{in}\ a\cdot b.
\]
\textsc{$\otimes$-Intro} gives
\[
\inferrule[$\otimes$-Intro]
  { \varnothing\mid x:\real{\E}\vdash x:\real{\E} \\
    \varnothing\mid y:\real{\E}\vdash y:\real{\E} \\
    \wf(x:\real{\E},y:\real{\E}) }
  { \varnothing\mid(x:\real{\E},y:\real{\E})\vdash(x,y):\real{\E}\otimes \real{\E} }.
\]
Unpacking this pair introduces the separated bindings
$(a:\real{\E},b:\real{\E})$, so \textsc{Mul-Ind} gives
\[
\varnothing\mid(a:\real{\E},b:\real{\E})\vdash a\cdot b:\real{\E}.
\]
Consequently, \textsc{$\otimes$-Elim} yields
\[
\varnothing\mid(x:\real{\E},y:\real{\E})\vdash
\mathsf{let}\ (a,b)=(x,y)\ \mathsf{in}\ a\cdot b:\real{\E}.
\]
\end{example}


\newpage


\section{Miscellaneous}
\begin{enumerate}
  \item In BI, contexts are tree-shaped structures called bunches. They distinguish between two ways of combining assymptions:
  \[
    \Gamma ::= x:\tau
    \mid \Gamma;\Gamma
    \mid \Gamma,\Gamma
  \]

  where:
  \begin{itemize}
    \item ; additive
    \item , multiplicative
  \end{itemize}

  \item Flat contexts vs bunched contexts 
  \item \textsc{IF} rule is not a problem because determinization does not eliminate the $G$ randomness, it retains it.
  
  \item let a = if z then x else y in
let b = if z then y else x in
a * b Currently, rejects the program.

however, you can rewrite it to the equivalent program:

if z then
  x * y
else
  y * x

\item In traditional BI presentations, the primitive additive typing rule usually uses the same context in both premises:
\[
\frac{
\Delta\mid\Gamma\vdash e_1:\real{\E}
\qquad
\Delta\mid\Gamma\vdash e_2:\real{\E}
}{
\Delta\mid\Gamma\vdash e_1+e_2:\real{\E}
}
\]

Your \(\Gamma_1;\Gamma_2\) formulation is an equivalent, convenient dependency-combining formulation, assuming the usual additive weakening and contraction principles. It makes context inference more explicit.
\end{enumerate}

\section{Related Work}
https://www.cambridge.org/core/journals/journal-of-functional-programming/article/on-bunched-typing/4AFEBD081F0F0AFFCF03C8415DDE53A4


\end{document}
